\chapter{Hemolytic Uremic Syndrome (HUS)}
\index{Hemolytic Uremic Syndrome (HUS)}
\label{sec:Hemolytic Uremic Syndrome (HUS)}

\section{Overview}
Hemolytic uremic syndrome (HUS) is a thrombotic microangiopathy characterized by the triad of microangiopathic hemolytic anemia, thrombocytopenia, and acute kidney injury. The most common form is typical (Shiga toxin--mediated) HUS, which predominantly affects children following infection with Shiga toxin--producing \emph{Escherichia coli} (STEC). Atypical HUS (aHUS) is a complement-mediated disorder with a distinct pathogenesis and management.

\section{Classification}
\begin{description}[font=\normalfont\bfseries,leftmargin=3em,labelwidth=2.5em]
  \item[Cause] STEC infection (typical HUS) or complement dysregulation (atypical HUS)
  \item[Organ System] Kidneys
  \item[Pathophysiology] Typical HUS: Shiga toxin binds to glomerular endothelial cells, inhibits protein synthesis, and triggers microvascular endothelial damage, platelet aggregation, and microthrombus formation. Atypical HUS: genetic or acquired defects in complement regulatory proteins (CFH, CFI, MCP) lead to uncontrolled alternative pathway activation and endothelial injury.
  \item[Duration] Acute with potential for chronic sequelae
  \item[Onset] Sudden (typically 3--10 days after diarrheal illness)
  \item[Mode of Transmission] Fecal-oral (typical HUS); not transmissible (atypical)
  \item[Clinical Course] Prodrome of diarrhea (often bloody), followed by pallor, oliguria, and signs of microangiopathic hemolytic anemia.
  \item[Severity] Moderate to severe
  \item[Extent of Disease] Systemic (hematologic, renal, and potentially neurologic involvement)
  \item[Age Group] Typical HUS: children under 5 years; atypical HUS: all ages
  \item[Epidemiology] Typical HUS: annual incidence 1--2 per 100,000 children; leading cause of acute kidney injury in children. Outbreaks associated with contaminated food (ground beef, leafy greens). Atypical HUS: rare (1--2 per million).
  \item[ICD-11 Code] 3A21.0 (typical), 3A21.1 (atypical)
\end{description}

\section{Causes}
Typical HUS is caused by infection with Shiga toxin--producing \emph{E. coli} (STEC), most commonly serotype O157:H7. Atypical HUS results from genetic mutations in complement regulatory genes (CFH, CFI, CFB, C3, MCP/THBD) or acquired autoantibodies against complement factor H.

\section{Risk Factors}
\begin{itemize}[noitemsep]
  \item Young age (under 5 years)
  \item Consumption of undercooked ground beef, unpasteurized dairy, or contaminated produce
  \item Exposure to contaminated water or infected individuals
  \item Antibiotic use during STEC infection (increases risk of HUS)
  \item Family history of aHUS or complement gene mutation
  \item Pregnancy or postpartum period (triggers aHUS)
\end{itemize}

\section{Signs and Symptoms}
\begin{itemize}[noitemsep]
  \item Bloody diarrhea (prodrome)
  \item Pallor and fatigue (from hemolytic anemia)
  \item Oliguria or anuria (acute kidney injury)
  \item Easy bruising or petechiae (thrombocytopenia)
  \item Neurologic symptoms (seizures, altered consciousness, stroke) in severe cases
  \item Difficulty performing daily activities
\end{itemize}

\section{Progression and Stages}
Typical HUS begins 3--10 days after onset of diarrheal illness. The acute phase lasts 1--3 weeks and is characterized by rapidly falling hemoglobin and platelet counts, with rising creatinine. Most children recover renal function with supportive care, but 20--50\% may require dialysis. Neurologic involvement indicates severe disease. Atypical HUS follows a relapsing or progressive course if untreated.

\section{Complications}
\begin{itemize}[noitemsep]
  \item Chronic kidney disease or end-stage renal disease
  \item Arterial hypertension (persistent after recovery)
  \item Neurologic sequelae (stroke, cognitive impairment, epilepsy)
  \item Cardiac involvement (cardiomyopathy, myocardial infarction)
  \item Reduced quality of life
  \item Emotional and psychological effects
\end{itemize}

\section{Diagnosis}
\begin{itemize}[noitemsep]
  \item Complete blood count showing anemia and thrombocytopenia
  \item Peripheral blood smear demonstrating schistocytes (microangiopathic hemolytic anemia)
  \item Elevated creatinine, LDH, and indirect bilirubin; low haptoglobin
  \item Lab tests (stool culture or PCR for STEC; complement studies and genetic testing for aHUS)
  \item Imaging tests (renal ultrasound to assess kidney size and echogenicity)
\end{itemize}

\section{Prevention}
\begin{itemize}[noitemsep]
  \item Cook ground beef to an internal temperature of 71°C (160°F)
  \item Practice hand hygiene and safe food handling
  \item Avoid unpasteurized dairy and untreated water
  \item Avoid antibiotics during STEC infection (may increase HUS risk)
  \item Regular check-ups
\end{itemize}

\section{Treatment Options}
\begin{itemize}[noitemsep]
  \item Supportive care: intravenous fluids, electrolyte management, blood transfusion
  \item Renal replacement therapy (dialysis) for acute kidney injury
  \item Eculizumab (complement C5 inhibitor) for atypical HUS
  \item Plasmapheresis/plasma exchange (for aHUS or severe neurologic involvement)
  \item Avoid platelet transfusion unless life-threatening bleeding
  \item Lifestyle changes
\end{itemize}

\section{Living with It}
\begin{itemize}[noitemsep]
  \item Follow treatment plan
  \item Regular appointments
  \item Adjust daily habits
  \item Support groups
  \item Keep learning
\end{itemize}

\section{Outlook}
In typical HUS, 90--95\% of children survive the acute episode, and most recover renal function fully. However, 20--30\% develop long-term renal sequelae (proteinuria, hypertension, reduced GFR). Atypical HUS has a poorer prognosis without treatment; eculizumab has dramatically improved outcomes, with most patients achieving hematologic and renal remission. Relapses occur in aHUS if treatment is discontinued.
